\section*{Vijay's Mysterious Laundry Counter}

At the University of Utah's \textbf{Kahlert School of Computing}, Professor
\textbf{Vijay} has ``modernized'' the \textbf{Hydrothermal Fabrication
Laboratory} (the very serious grad-student laundromat).

He installs a big digital display on the wall that shows a single number:

\begin{quote}
the \textbf{current number of finished laundry loads} for the day.
\end{quote}

Several student assistants help run the lab. During the day, they:
\begin{itemize}[nosep]
    \item Update the number on the display, and
    \item Check the number to fill out various forms and reports.
\end{itemize}

To ``avoid crowding around the display,'' Vijay gives each assistant a quirky
set of rules for how they record and apply changes.

\subsection*{How the Assistants Work}

There is:
\begin{itemize}[nosep]
    \item One \textbf{shared display} showing a single integer \(X\) (starting
    from some initial value).
    \item \(T\) assistants, numbered \(1\) through \(T\).
\end{itemize}

Each assistant keeps a \textbf{personal clipboard} where they write down planned
updates \textbf{before} walking over to the display.

Each assistant's program is a list of instructions of three kinds:
\begin{enumerate}[label=\arabic*., leftmargin=*]
    \item \texttt{WRITE v}
    \begin{itemize}
        \item The assistant writes down a note on their clipboard saying ``the
        display should show \texttt{v}''.
        \item This does \textbf{not} change the shared display yet.
        \item The assistant keeps all such notes in order on their clipboard.
    \end{itemize}

    \item \texttt{POST}
    \begin{itemize}
        \item The assistant walks over to the display and applies \textbf{all of
        their notes}, in order.
        \item For each note with value \(v\), they set the display to \(v\).
        \item After applying all notes, they clear their clipboard (no notes
        left).
    \end{itemize}

    \item \texttt{CHECK}
    \begin{itemize}
        \item The assistant wants to read what the ``current'' number of
        finished loads is.
        \item However, they trust their own notes more than the distant display:
        \begin{itemize}
            \item If their clipboard has \textbf{at least one note}, they read
            the value from the \textbf{most recent} note on their clipboard (the
            last \texttt{WRITE} they did that hasn't been cleared by a
            \texttt{POST}).
            \item If their clipboard is \textbf{empty}, they read the value from
            the shared display \(X\).
        \end{itemize}
        \item Whenever a \texttt{CHECK} happens, we want to \textbf{record} the
        number they saw.
    \end{itemize}
\end{enumerate}

Different assistants may have different lists of instructions.

\subsection*{The Daily Schedule}

You are also given a \textbf{schedule} for the day:
\begin{itemize}
    \item The schedule is a sequence of assistant IDs.
    \item At step \(k\) of the day, the assistant with ID \(s_k\) performs
    \textbf{their next instruction}, if they still have one.
    \item If an assistant is scheduled but already finished all their
    instructions, they simply do nothing at that step.
\end{itemize}

All assistants share the same display \(X\), but each one has their own
clipboard of notes.

Your job is to simulate this entire process and find out exactly what numbers
were seen during all the \texttt{CHECK} operations.

\subsection*{Input Format}

\begin{enumerate}[label=\arabic*., leftmargin=*]
    \item First line: three integers
\begin{verbatim}
T P X0
\end{verbatim}
    \begin{itemize}
        \item \(T\) --- number of assistants (\(1 \le T \le 10\))
        \item \(P\) --- length of the schedule (\(1 \le P \le 100000\))
        \item \(X_0\) --- initial value shown on the display (\(0 \le X_0 \le
        10^9\))
    \end{itemize}

    \item For each assistant \(i\) from \(1\) to \(T\):
    \begin{itemize}
        \item One line: integer \(n_i\) --- the number of instructions for
        assistant \(i\)\\
        (\(0 \le n_i \le 100000\), and the sum of all \(n_i\) over all
        assistants \(\le 100000\))
        \item Then \(n_i\) lines, each containing one instruction, in one of
        these forms:
        \begin{itemize}
            \item \texttt{WRITE v} --- where \(0 \le v \le 10^9\)
            \item \texttt{POST}
            \item \texttt{CHECK}
        \end{itemize}
    \end{itemize}

    \item Final line: \(P\) integers
\begin{verbatim}
s_1 s_2 ... s_P
\end{verbatim}
    Each \(s_k\) is in the range \(1..T\) and indicates that at step \(k\),
    assistant \(s_k\) takes their next instruction (if any).
\end{enumerate}

If an assistant has no instructions left when chosen in the schedule, they do
nothing for that step.

You may assume that no invalid situations (like negative values) occur.

\subsection*{Output Format}

For \textbf{every} \texttt{CHECK} instruction executed during the schedule
(across all assistants), print the value observed on its own line, in the exact
order the \texttt{CHECK} operations happen.

If there are \(K\) \texttt{CHECK} operations executed in total, your output
should have exactly \(K\) lines.

\subsection*{Example}

\textbf{Input}
\begin{verbatim}
2 7 0
3
WRITE 5
CHECK
POST
3
WRITE 10
CHECK
POST
1 2 1 2 1 2 2
\end{verbatim}

\textbf{Explanation}

There are 2 assistants. Initial display value: \(X = 0\).

Assistant 1's instructions:
\begin{enumerate}[nosep]
    \item \texttt{WRITE 5}
    \item \texttt{CHECK}
    \item \texttt{POST}
\end{enumerate}

Assistant 2's instructions:
\begin{enumerate}[nosep]
    \item \texttt{WRITE 10}
    \item \texttt{CHECK}
    \item \texttt{POST}
\end{enumerate}

Schedule: \(1\ 2\ 1\ 2\ 1\ 2\ 2\).

\begin{itemize}
    \item Step 1: assistant 1 \(\rightarrow\) \texttt{WRITE 5}
    \begin{itemize}
        \item Clipboard1: \([5]\), Clipboard2: \([]\)
        \item Display \(X = 0\)
    \end{itemize}

    \item Step 2: assistant 2 \(\rightarrow\) \texttt{WRITE 10}
    \begin{itemize}
        \item Clipboard1: \([5]\), Clipboard2: \([10]\)
        \item Display \(X = 0\)
    \end{itemize}

    \item Step 3: assistant 1 \(\rightarrow\) \texttt{CHECK}
    \begin{itemize}
        \item Clipboard1 not empty \(\Rightarrow\) sees last note: \(5\)
        \item Output: \(5\)
    \end{itemize}

    \item Step 4: assistant 2 \(\rightarrow\) \texttt{CHECK}
    \begin{itemize}
        \item Clipboard2 not empty \(\Rightarrow\) sees last note: \(10\)
        \item Output: \(10\)
    \end{itemize}

    \item Step 5: assistant 1 \(\rightarrow\) \texttt{POST}
    \begin{itemize}
        \item Applies notes in order: sets \(X = 5\)
        \item Clipboard1 becomes empty
    \end{itemize}

    \item Step 6: assistant 2 \(\rightarrow\) \texttt{POST}
    \begin{itemize}
        \item Applies notes in order: sets \(X = 10\)
        \item Clipboard2 becomes empty
    \end{itemize}

    \item Step 7: assistant 2 \(\rightarrow\) no instructions left, does
    nothing.
\end{itemize}

Final \texttt{CHECK} results, in order:
\begin{verbatim}
5
10
\end{verbatim}

Professor Vijay happily announces that this proves his ``laundry coordination
rules'' are perfectly reasonable.

The students are\ldots{} not entirely convinced.
